% マルチソース構成のサンプル。
% 1 ファイル = 1 モジュールで、古いスライドは古いマクロライブラリのまま、
% 新しいスライドは新しいライブラリのまま、1 つの文書に共存する。


\title{レイテンシ改善の記録}
\maketitle

% 呼び出し先の既定値のまま取り込む。


\begin{frame}{2023 年の測定結果}
{
\color{blue}
当時のベースラインは p99 で 420 ms でした。
}

\end{frame}

% 呼び出し先のフラグをリテラルで上書きする。


\begin{frame}{現在の測定結果}
{
\color{orange!80!black}\bfseries
いまのベースラインは p99 で 150 ms です。
}

測定環境: 8 vCPU / 16 GiB、同時実行数 64。
\end{frame}

% 呼び出し側のフラグを明示的に転送する。同名でも暗黙には継承されない。


\begin{frame}{2023 年の測定結果}
{
\color{blue}
当時のベースラインは p99 で 420 ms でした。
}

測定環境: 4 vCPU / 8 GiB、同時実行数 64。
\end{frame}

% 条件が false なら、そのファイルは開かれもしない。
